\chapter{Konfigurasi Awal}

\section{Perintah Setup Otomatis}

PakAI menyediakan perintah \cmd{pakai setup} yang mendeteksi penyedia yang terinstal,
menulis unit systemd, dan menampilkan cuplikan konfigurasi yang siap digunakan.

\begin{lstlisting}[style=terminal]
$ pakai setup
\end{lstlisting}

\begin{figure}[h]
  \centering
  \includegraphics[width=0.88\linewidth]{pakai-setup}
  \caption{Output \texttt{pakai setup}: deteksi penyedia dan konfigurasi otomatis}
\end{figure}

\subsection{Yang Dilakukan oleh \texttt{pakai setup}}

\begin{enumerate}
  \item \textbf{Deteksi penyedia}: memeriksa jalur berikut:
    \begin{itemize}
      \item Claude: \texttt{\textasciitilde/.claude/stats-cache.json}
      \item OpenCode Go: \texttt{\textasciitilde/.local/share/opencode/opencode-stable.db}
            atau \texttt{opencode.db}
      \item OpenAI Codex: \texttt{\textasciitilde/.codex/auth.json}
    \end{itemize}
  \item \textbf{Memulai daemon} jika belum berjalan.
  \item \textbf{Membuat unit systemd} di\\
        \texttt{\textasciitilde/.config/systemd/user/pakai.service}.
  \item \textbf{Menampilkan cuplikan konfigurasi} untuk tmux dan Waybar.
\end{enumerate}

\section{Memulai Daemon}

Daemon adalah proses latar belakang yang mengambil dan menyimpan data penggunaan.
Tanpa daemon, perintah-perintah PakAI tidak akan menampilkan data.

\subsection{Memulai daemon secara manual}

\begin{lstlisting}[style=terminal]
$ pakai daemon start
Daemon started (port 7731)
\end{lstlisting}

\subsection{Memeriksa status daemon}

\begin{lstlisting}[style=terminal]
$ pakai daemon status
\end{lstlisting}

\begin{figure}[h]
  \centering
  \includegraphics[width=0.72\linewidth]{pakai-daemon-status}
  \caption{Output \texttt{pakai daemon status}}
\end{figure}

\subsection{Menghentikan daemon}

\begin{lstlisting}[style=terminal]
$ pakai daemon stop
\end{lstlisting}

\section{Konfigurasi Autostart dengan Systemd}

Agar daemon berjalan otomatis saat login, aktifkan unit systemd yang dibuat oleh
\cmd{pakai setup}:

\begin{lstlisting}[style=terminal]
$ systemctl --user enable --now pakai
$ systemctl --user status pakai
\end{lstlisting}

Isi unit systemd yang dibuat secara otomatis:

\begin{lstlisting}[style=toml]
[Unit]
Description=PakAI usage tracker daemon
After=default.target

[Service]
ExecStart=/home/user/.local/bin/pakai daemon start --foreground
Restart=on-failure
RestartSec=5

[Install]
WantedBy=default.target
\end{lstlisting}

\begin{tipbox}
  Gunakan \cmd{journalctl --user -u pakai -f} untuk melihat log daemon secara
  langsung (\textit{real-time}).
\end{tipbox}

\section{Autentikasi Penyedia}

\subsection{Claude}

Claude menggunakan OAuth yang dikelola oleh Claude Code CLI. Jalankan:

\begin{lstlisting}[style=terminal]
$ claude login
\end{lstlisting}

Kredensial disimpan di \texttt{\textasciitilde/.claude/.credentials.json}.

\subsection{OpenAI Codex}

\begin{lstlisting}[style=terminal]
$ codex login
\end{lstlisting}

Kredensial disimpan di \texttt{\textasciitilde/.codex/auth.json}.

\subsection{OpenCode Go}

OpenCode Go tidak memerlukan login tambahan. PakAI membaca langsung dari database
SQLite yang dibuat oleh OpenCode. Pastikan OpenCode sudah pernah digunakan minimal
satu kali sehingga database terbentuk di:

\begin{lstlisting}[style=terminal]
~/.local/share/opencode/opencode-stable.db
\end{lstlisting}
